\chapter*{ABSTRACT}
\phantomsection
\addcontentsline{toc}{chapter}{ABSTRACT}

\bigskip

Ambient air pollution is a major environmental and public-health concern in urban areas where pollutant concentration varies across space, time, and season. Kathmandu Valley is affected by road traffic, dust resuspension, construction activities, brick kilns, seasonal forest-fire episodes, valley topography, and changing meteorological conditions, making short-term PM$_{2.5}$ forecasting important for air quality interpretation and public awareness. Traditional station-wise forecasting models often ignore wind-driven pollutant movement between monitoring stations and usually provide numerical outputs without clearly explaining why a forecast is high or low. This proposal presents an explainable wind-aware Spatio-Temporal Graph Neural Network (ST-GNN) framework for short-term PM$_{2.5}$ forecasting, interpretation, and visualization in Kathmandu Valley. Monitoring stations are represented as graph nodes, while dynamic edges are constructed using geographical distance, historical PM$_{2.5}$ correlation, wind speed, and wind-direction alignment. The primary forecasting model is a Graph Attention Network with Gated Recurrent Unit (GAT-GRU), where GAT captures spatial influence between stations and GRU captures temporal patterns from historical observations. The forecasting task is formulated as a regression problem for hourly PM$_{2.5}$ prediction at a 6-hour horizon, while Air Quality Index categories and pollution-peak labels are derived from predicted concentrations to support interpretation and event-level evaluation. The system includes timestamp alignment, missing-value handling, sensor anomaly detection, outlier treatment, normalization, and explainability through attention-based station influence, feature attribution, and temporal sensitivity analysis. The final output will be an interactive dashboard showing station-level PM$_{2.5}$ forecasts, AQI risk classes, historical-versus-predicted trends, spatial pollution maps, virtual-node estimates for unmonitored locations, and feature-importance visualizations. The project is intended as a research prototype and decision-support tool for air quality interpretation and visualization, not as a replacement for official air-quality monitoring or warning systems.

\vspace{1em}
\noindent\textit{\textbf{Keywords}---Air quality forecasting, PM$_{2.5}$, wind-aware graph, spatio-temporal graph neural network, GAT-GRU, explainable AI, AQI, Kathmandu Valley.}
